% Adaptive Intelligent Disaster Response Agent (AIDRA) - IEEE Conference LaTeX
% Save this file as report.tex and compile with pdflatex / latexmk in Overleaf.
\documentclass[conference]{IEEEtran}
\IEEEoverridecommandlockouts
\usepackage{times}
\usepackage{graphicx}
\usepackage{amsmath,amssymb}
\usepackage{booktabs}
\usepackage{caption}
\usepackage{subcaption}
\usepackage{balance}
\usepackage{url}
\usepackage{hyperref}
\usepackage{siunitx}
\usepackage{algorithm}
\usepackage{algpseudocode}
\usepackage{microtype}

% Slight compression to help meet 6-page limit
\setlength{\textfloatsep}{7pt}
\setlength{\floatsep}{7pt}
\setlength{\intextsep}{7pt}
\renewcommand{\baselinestretch}{0.98}

\begin{document}

\title{Adaptive Intelligent Disaster Response Agent (AIDRA):\\
A Hybrid System for Dynamic Rescue Planning and Resource Allocation}

\author{\IEEEauthorblockN{Group Member 1, Group Member 2}
\IEEEauthorblockA{Course: ARTIFICIAL INTELLIGENCE (AIC 201)\\
Instructor: Dr. Arshad Farhad, \texttt{arshadfarhad.buic@bahria.edu.pk}\\
Assignment: Assignment 3 \& Assignment 4 (CCP-Based)}}

\maketitle

\begin{abstract}
This report documents the design, implementation, and evaluation of the
Adaptive Intelligent Disaster Response Agent (AIDRA). AIDRA integrates search-based
planning, constraint satisfaction, local search, machine learning risk estimation,
and uncertainty handling to operate under conflicting objectives (rescue time vs.
risk exposure, victim prioritization vs. throughput) and dynamic environmental events
(road blockages, spreading hazards). The document provides the system architecture,
algorithms, experimental protocol, KPI computation, decision logs, and placeholders
for artifacts (code repository and demo video).
\end{abstract}

\begin{IEEEkeywords}
disaster response, planning, A*, CSP, fuzzy logic, machine learning, dynamic replanning
\end{IEEEkeywords}

\section{Introduction}
Disaster response planning requires fast, robust decisions under uncertainty.
We present AIDRA, a modular hybrid agent that combines classical search, CSP-based
resource allocation, local search optimization, supervised ML risk estimators,
and fuzzy reasoning to make explainable decisions in a dynamic simulated urban
disaster (the provided CCP). The implementation and runtime traceability are
intended to demonstrate trade-offs between conflicting objectives and to enable
comparative evaluation.

\section{CCP Summary and Objectives}
Concise re-statement of the CCP baseline:
\begin{itemize}
  \item Grid-based urban map with high-risk zones and potentially blockable roads.
  \item Resources: 2 ambulances, 1 rescue team, 10 medical kits.
  \item Victims: 5 total (2 critical, 2 moderate, 1 minor).
  \item Required outputs per run: prioritized rescue order, per-trip route with
  explicit time-vs-risk trade-off, resource-allocation plan respecting hard
  constraints, survival/risk estimates at rescue time, replanning decision log,
  and comparative performance report across strategies.
\end{itemize}

\section{System Architecture}
Figure~\ref{fig:arch} shows the high-level components and dataflow.

\begin{figure}[!t]
  \centering
  % Insert a high-quality PDF of the architecture into figs/architecture.pdf
  \includegraphics[width=0.95\linewidth]{figs/architecture.pdf}
  \caption{AIDRA high-level architecture (placeholder).}
  \label{fig:arch}
\end{figure}

Implementation mapping (files referenced in repository):
\begin{itemize}
  \item Search and heuristics: \texttt{src/engine/search.ts}, \texttt{src/engine/heuristics.ts}
  \item Grid and environment generation: \texttt{src/engine/gridGenerator.ts}
  \item Constraint solver: \texttt{src/engine/csp.ts}
  \item Local search: \texttt{src/engine/localSearch.ts}
  \item ML risk pipeline: \texttt{src/engine/mlRiskPipeline.ts}
  \documentclass[conference]{IEEEtran}
  \IEEEoverridecommandlockouts
  \usepackage{times}
  \usepackage{graphicx}
  \usepackage{amsmath,amssymb}
  \usepackage{booktabs}
  \usepackage{caption}
  \usepackage{subcaption}
  \usepackage{balance}
  \usepackage{url}
  \usepackage{hyperref}
  \usepackage{siunitx}
  \usepackage{algorithm}
  \usepackage{algpseudocode}
  \usepackage{microtype}
  \usepackage{listings}
  \usepackage{longtable}

  % Slight compression to help meet 6-page limit but keep readable
  \setlength{\textfloatsep}{6pt}
  \setlength{\floatsep}{6pt}
  \setlength{\intextsep}{6pt}
  \renewcommand{\baselinestretch}{0.985}

  \begin{document}

  	itle{Adaptive Intelligent Disaster Response Agent (AIDRA):\\
  A Hybrid System for Dynamic Rescue Planning and Resource Allocation}

  \author{\IEEEauthorblockN{Group Member 1, Group Member 2}
  \IEEEauthorblockA{Course: ARTIFICIAL INTELLIGENCE (AIC 201)\\
  Instructor: Dr. Arshad Farhad, \texttt{arshadfarhad.buic@bahria.edu.pk}\\
  Assignment: Assignment 3 \& Assignment 4 (CCP-Based)}}

  \maketitle

  \begin{abstract}
  This report presents the design, implementation, and empirical evaluation of
  the Adaptive Intelligent Disaster Response Agent (AIDRA). AIDRA integrates
  classical search and planning, constraint satisfaction, local search optimization,
  machine learning-based risk estimation, and fuzzy uncertainty handling to
  address multi-objective rescue problems under dynamic hazards. The system is
  designed for the CCP baseline: five victims on a grid, two ambulances, one
  rescue team, and changing environment events (blocked roads, spreading fires).
  This document includes formal problem definitions, algorithms, implementation
  mapping to the codebase, a detailed experimental protocol, KPI computation
  formulas, example decision logs, and placeholders for artifact links.
  \end{abstract}

  \begin{IEEEkeywords}
  disaster response, planning, A*, CSP, fuzzy logic, machine learning, dynamic replanning
  \end{IEEEkeywords}

  \section{Introduction}
  Disaster response involves making critical trade-offs under uncertainty and
  resource constraints. The CCP requires a system-level solution where multiple
  AI techniques interact to generate, explain, and revise plans in real time.
  This work focuses on an integrated pipeline that: (1) models the environment and
  objectives, (2) generates route and assignment plans using search and CSP,
  (3) estimates risk with ML models to inform decisions, and (4) handles uncertainty
  via fuzzy reasoning and explicit logging for traceability.

  We aim to demonstrate how different algorithmic choices affect KPIs such as
  victims saved, average rescue time, and cumulative risk exposure, and to
  provide reproducible artifacts (code and demonstration video) as required.

  \section{CCP Formal Problem Definition}
  We formalize the CCP as a dynamic, partially-observable multi-agent planning
  problem with hard resource constraints.

  \subsection{Environment and Entities}
  - Grid: a 2D grid $G=(V,E)$ where $V$ are cells and $E$ edges representing
    traversable connections. Each cell $v\in V$ has attributes: risk level
    $r(v)\in[0,1]$, traversability $t(v)\in\{0,1\}$, and possibly a victim.
  - Agents: two ambulances $A_1,A_2$ and one rescue team $R$.
  - Victims: set $\mathcal{V} = \{v_1,\dots,v_5\}$ with severity levels
    $s(v_i)\in\{\text{critical, moderate, minor}\}$.
  - Medical centers: two fixed cells $M_1,M_2$ with differing distances.

  \subsection{State, Actions, Observations}
  State $S_t$ at time $t$ includes positions of ambulances and rescue team,
  victim states (location, health, rescued flag), road states (open/blocked),
  and kit inventory $K_t$. Actions $a\in A$ include movement along edges,
  pick-up, drop-off, and reassign commands. Observations $o_t$ include explicit
  event messages: blockages, hazard intensity updates, and discovered victims.

  \subsection{Objectives}
  We treat the problem as multi-objective. Primary objectives and formal metrics:
  \begin{align}
    &\text{VictimsSaved} = \sum_{i=1}^{5} \mathbf{1}\{\text{victim }i\text{ reached medical center}\}\\
    &\text{AvgRescueTime} = \frac{1}{\max(1,\text{VictimsSaved})} \sum_{i:\text{rescued}} T_{rescue}(i)\\
    &\text{RiskExposure} = \sum_{agents}\sum_{t} r(cell(agent,t))\cdot\Delta t
  \end{align}

  Agents must satisfy the hard constraint: at any time, an ambulance may carry at
  most two victims.

  \section{System Architecture (Detailed)}
  Figure~\ref{fig:arch} provides the high-level module layout. Each module is
  described with responsibilities and interfaces.

  \begin{figure}[!t]
    \centering
    \includegraphics[width=0.98\linewidth]{figs/architecture.pdf}
    \caption{AIDRA system architecture (placeholder PDF).}
    \label{fig:arch}
  \end{figure}

  \subsection{Module Responsibilities}
  - \textbf{Simulator / Environment}: manages grid, injects events, updates road
    and risk states (\texttt{src/engine/simulationEngine.ts}).
  - \textbf{Search \& Planner}: route computation and replanning (BFS, DFS,
    Greedy Best-First, A*; see \texttt{src/engine/search.ts}).
  - \textbf{CSP Solver}: assignment of victims to ambulances and scheduling
    (backtracking with MRV/degree heuristics; \texttt{src/engine/csp.ts}).
  - \textbf{Local Search}: schedule/post-optimization operators (\texttt{src/engine/localSearch.ts}).
  - \textbf{ML Risk Pipeline}: feature extraction, model training, inference
    (\texttt{src/engine/mlRiskPipeline.ts}).
  - \textbf{Uncertainty Handler}: fuzzy rule base and probability estimators
    (\texttt{src/engine/fuzzyLogic.ts}).
  - \textbf{Logger / Visualizer}: decision logs, search traces, and export of
    KPIs/figures (e.g., \texttt{src/data/placeholder.ts}).

  \section{Search Algorithms and Heuristics}
  We implement and evaluate the following planners with complexity notes and
  practical considerations.

  \subsection{Uninformed Search (BFS/DFS)}
  Used for correctness baseline and small instance comparisons. BFS guarantees
  shortest-path in unweighted grids; DFS provides quick exploration in constrained
  settings but no optimality.

  \subsection{Greedy Best-First}
  Expands nodes by estimated proximity to goal using heuristic $h(n)$ (fast but
  not optimal). Suitable for time-prioritization when risk term is small.

  \subsection{A* (Admissible Heuristic)}
  A* is used with a combined heuristic:
  \begin{equation}
    h(n) = \alpha \cdot \mathrm{dist}_{time}(n,goal) + \beta \cdot \mathrm{risk}_{est}(n,goal)
  \end{equation}
  Parameters $\alpha,\beta\ge 0$ control the trade-off. We ensure admissibility by
  keeping $\beta$ small relative to an admissible lower bound on added risk cost
  or by using a heuristic that underestimates true cost.

  \subsection{Search Traceability}
  Each search run logs: nodes expanded, frontier size over time, g-values, and
  selected path. These traces are exported for visualization (search trace PDF).

  \section{Constraint Satisfaction Formulation}
  We present the CSP formalization used to allocate victims and schedule trips.

  \subsection{Variables and Domains}
  - For each ambulance trip slot $s$ and ambulance $a$, variable $X_{a,s}$ denotes
    the subset of victims assigned (domain = subsets of remaining victims with
    size $\le 2$).
  - For rescue team timeslots, $Y_t$ denotes which location is serviced.

  \subsection{Constraints}
  - Capacity: $|X_{a,s}|\le 2$.
  - Exclusivity: a victim can be assigned to at most one ambulance trip.
  - Temporal feasibility: assigned trips must respect travel-time windows given
    current road states.

  \subsection{Solver Enhancements}
  We use MRV to pick variables with smallest remaining domain, degree heuristic to
  break ties, and forward checking to prune inconsistent domains early. We record
  backtrack counts and runtime to compare heuristic effectiveness.

  \section{Local Search Details}
  Local search is applied to candidate assignment solutions to reduce objective
  value post-CSP. Operators include swap two victims between ambulances, move one
  victim to a different trip, and reverse route segment. For Simulated Annealing
  we use an exponential cooling schedule $T_k=T_0\cdot\gamma^k$ with $\gamma\in(0,1)$.

  \section{Machine Learning Risk Pipeline}
  This section describes dataset synthesis, feature set, models, evaluation, and
  integration into decision-making.

  \subsection{Synthetic Dataset Generation}
  When field data is unavailable, we synthesize labeled examples per grid cell and
  victim by sampling event times, hazard intensities, victim severities, and
  simulating short trajectories. Each example includes features listed below and
  a binary survival label under a fixed rescue delay.

  \subsection{Feature Set}
  \begin{itemize}
    \item Victim severity (one-hot encoded)
    \item Distance to nearest medical center (Euclidean / grid hops)
    \item Estimated travel time from nearest ambulance
    \item Local hazard intensity (continuous)
    \item Number of expected blockages on route (probabilistic estimate)
    \item Time since incident reported
  \end{itemize}

  \subsection{Models and Training}
  We implement and compare k-NN, Naive Bayes, and a small MLP (1 hidden layer).
  Training uses stratified k-fold cross-validation (k=5) and hyperparameter
  search (grid search) for MLP hidden units and k in k-NN. Evaluation metrics are
  reported per fold and aggregated.

  \subsection{Integration}
  Model outputs $p_{survival}(v)$ are converted into risk contributions used in
  heuristic cost: $r_{est} = 1 - p_{survival}$ (scaled). Policies may prioritize
  victims with lower survival probability to maximize expected survivors.

  \section{Uncertainty and Fuzzy Logic}
  We implement a fuzzy rule base to combine noisy indicators into an interpretable
  risk label. Membership functions are triangular/trapezoidal; example variable
  "hazard intensity" uses Low/Medium/High sets.

  Example rule (informal):
  \begin{quote}
  IF hazardIntensity IS High AND proximity IS Near THEN risk IS VeryHigh.
  \end{quote}

  Defuzzified numeric risk is used both as a planner input and as an explanatory
  term in the decision log.

  \section{Dynamic Replanning Policy}
  We implement an event-driven replanning manager with the following logic:
  \begin{enumerate}
    \item On event, compute impact on active plans (affected trips, delay estimates).
    \item If impact exceeds threshold (e.g., added delay > $\delta$ or route
      blockage on critical trip), trigger full replanning: recompute CSP and route
      plans using current world-state.
    \item Otherwise attempt local repair: reroute affected ambulance only.
  \end{enumerate}

  Every decision produces a log entry recording: timestamp, event type, previous
  plan summary, new plan summary, objective weights used, and justification.

  \section{Experimental Design and Protocol}
  This section details how to run controlled experiments and how to vary
  parameters to compare algorithmic choices.

  \subsection{Baseline Scenario}
  Use the CCP baseline: grid size 16x16, victims at predefined coordinates, two
  medical centers, 2 ambulances, 1 rescue team, 10 kits. Use fixed random seeds
  for event injection to allow reproducibility.

  \subsection{Parameter Sweeps}
  Key parameter dimensions:
  \begin{itemize}
    \item Objective weight sweep: $\alpha\in\{0.2,0.5,1.0\}$, $\beta\in\{0.2,0.5,1.0\}$.
    \item Planner: {BFS, DFS, Greedy, A*}.
    \item CSP heuristics: {none, MRV\&degree\&forward-checking}.
    \item ML models: {k-NN, MLP}.
  \end{itemize}

  For each experimental cell, run $N=30$ stochastic trials and report mean and
  95\% confidence intervals for each KPI.

  \section{KPIs: Definitions and Computation}
  We provide formulas used to compute each KPI in the repository to ensure
  consistent evaluation.

  \subsection{Victims Saved and Failure Modes}
  VictimsSaved is the integer count of victims reaching any medical center before
  their clinically modeled survival horizon. Report both absolute counts and
  percentage.

  \subsection{Average Rescue Time}
  Average over rescued victims; for non-rescued victims, time is counted until
  simulation end and reported separately.

  \subsection{Path Optimality Ratio}
  When best-known cost is available (small instances), compute
  \[ \text{OptimalityRatio} = \frac{cost_{found}}{cost_{best}}. \]

  \subsection{Resource Utilization Rate}
  Computed as fraction of time ambulances and team spend performing useful work
  (traveling with victim or active rescue) over total simulation time.

  \subsection{Risk Exposure Score}
  Sum of risk-weighted time spent by all agents in risk zones, optionally weighted
  by victim severity.

  \section{Example Scenario (Illustrative)}
  Below is an illustrative baseline instance you can run and use as an example
  in the report. Replace with real experiment outputs when available.

  \begin{itemize}
    \item Grid: 16x16
    \item Base (ambulance origin): (1,1)
    \item Medical centers: M1=(16,1), M2=(16,16)
    \item Victim positions and severities: V1=(4,5,critical), V2=(5,12,critical),
      V3=(10,8,moderate), V4=(12,3,moderate), V5=(8,14,minor)
    \item Event: road blockage at edge ((6,5)-(6,6)) at t=30s
  \end{itemize}

  Sample (placeholder) outcomes below should be replaced with real run data:
  \begin{longtable}{p{0.32\linewidth}p{0.64\linewidth}}
  	oprule
  Outcome & Placeholder value \\\midrule
  Victims saved & 4/5 \\
  Avg rescue time & 235 s \\
  Risk exposure & 12.4 (risk-units) \\
  Planner used & A* with $\alpha=1.0,\beta=0.5$ \\
  ML model & MLP (accuracy 0.86) \\
  \bottomrule
  \end{longtable}

  \section{Decision Log Schema}
  We recommend this JSON schema for each log entry (exported to CSV/JSON):
  \begin{lstlisting}
  {
    "timestamp": "2026-05-05T00:01:34Z",
    "event": "road_blockage",
    "location": [6,5],
    "affected_plans": ["Amb1_trip2"],
    "previous_plan_summary": "Amb1: V1->M1 via path P1",
    "new_plan_summary": "Amb1: V1->M1 via path P2",
    "objective_weights": {"alpha":1.0,"beta":0.5},
    "justification": "Blocked edge on P1; alternative avoids high-risk cell"
  }
  \end{lstlisting}

  \section{Implementation Notes and File Map}
  Key files and their purposes (update if your code differs):
  \begin{itemize}
    \item \texttt{src/engine/search.ts} -- search algorithms (BFS/DFS/Greedy/A*).
    \item \texttt{src/engine/heuristics.ts} -- heuristic functions and weight sweep harness.
    \item \texttt{src/engine/gridGenerator.ts} -- grid and scenario builders.
    \item \texttt{src/engine/csp.ts} -- CSP solver with heuristics and instrumentation.
    \item \texttt{src/engine/localSearch.ts} -- SA/Hill-climbing routines.
    \item \texttt{src/engine/mlRiskPipeline.ts} -- feature pipeline, model training, scoring.
    \item \texttt{src/engine/fuzzyLogic.ts} -- fuzzy rule base and defuzzification.
    \item \texttt{src/engine/simulationEngine.ts} -- event injection and time loop.
    \item \texttt{src/data/placeholder.ts} -- example scenarios and UI placeholders.
  \end{itemize}

  \section{Reproducibility Checklist}
  Include these artifacts in your GitHub repository prior to submission:
  \begin{itemize}
    \item Scenario JSONs and seeds (to reproduce experiments).
    \item Scripts to run experiments and aggregate KPIs (e.g., `scripts/run_all.sh`).
    \item Raw CSV logs, aggregated KPI CSV, and figure-generating scripts.
    \item `README.md` with exact commands, Node version, and dependency list.
    \item Demonstration video link and notes on what is shown.
  \end{itemize}

  \section{Guidelines for Figures and Tables}
  Use vector PDF for diagrams and high-resolution PNG for raster plots. Keep
  captions concise and reference files in `figs/` directory. All figures must be
  embedded as PDF in the final submission to preserve clarity.

  \section{Discussion, Limitations, and Future Work}
  We discuss effects of noisy sensors, simplified hazard dynamics, and the gap
  between synthetic and field data. Future work includes richer ML using real
  incident data, reinforcement-learning-based policies for continuous control,
  and distributed multi-ambulance coordination.

  \section{Conclusion}
  This extended LaTeX template provides the textual content and structure needed
  to produce an IEEE-format 6-page report for the AIDRA CCP assignment. Replace
  placeholders, insert experimental results and high-quality figures, and verify
  that the final PDF adheres to the 6-page limit.

  \section*{Acknowledgments}
  Instructor: Dr. Arshad Farhad.

  \balance
  \bibliographystyle{IEEEtran}
  \bibliography{refs}

  % ------------------------- Appendices (extended content) -------------------------
  \appendix
  \section{Detailed Algorithms and Pseudocode}
  This appendix contains detailed pseudocode for core algorithms used in AIDRA.

  \subsection{Backtracking CSP Solver with MRV and Forward Checking}
  The CSP solver is central to resource allocation. The pseudocode below mirrors
  the implementation in `src/engine/csp.ts` and documents instrumentation points
  for backtracks and domain pruning.

  \begin{algorithm}
  \caption{CSP Backtracking with MRV + Forward Checking}
  \begin{algorithmic}[1]
  \Procedure{CSP-Backtrack}{$Vars,Domains,Constraints$}
    \If{all variables assigned} \State \textbf{return} solution \EndIf
    \State var $\leftarrow$ SelectVariableMRV($Vars,Domains$)
    \ForAll{value $\in$ OrderDomain(var)}
      \If{Consistent(var,value,Constraints)}
        \State Assign(var,value)
        \State savedDomains $\leftarrow$ copy(Domains)
        \State Domains $\leftarrow$ ForwardCheck(var,value,Domains,Constraints)
        \If{no domain empty}
          \State result $\leftarrow$ CSP-Backtrack($Vars,Domains,Constraints$)
          \If{result successful} \State \textbf{return} result \EndIf
        \EndIf
        \State Domains $\leftarrow$ savedDomains
        \State Unassign(var)
        \State increment(backtrackCounter)
      \EndIf
    \EndFor
    \State \textbf{return} failure
  \EndProcedure
  \end{algorithmic}
  \end{algorithm}

  Notes: MRV chooses the variable with smallest domain; tie-breaking uses degree
  heuristic. Forward checking prunes neighbors' domains; every domain prune is
  logged for later analysis.

  \subsection{A* Implementation Notes}
  We implement A* with an adaptable heuristic function. To support trade-off
  analysis we parameterize the heuristic as a linear combination and expose
  weights to the experiment harness. Pseudocode is standard (see Section V), with
  instrumentation hooks to record nodes expanded and frontier sizes.

  \subsection{Simulated Annealing for Schedule Optimization}
  Key parameters used in experiments:
  \begin{itemize}
    \item Initial temperature $T_0 = 1.0$.
    \item Cooling factor $\gamma = 0.995$.
    \item Iteration cap $I_{max} = 10000$.
  \end{itemize}

  Pseudocode (high-level):
  \begin{algorithm}
  \caption{Simulated Annealing (schedule optimization)}
  \begin{algorithmic}[1]
  \Procedure{SA}{$s_0$}
    \State $s \leftarrow s_0$, $T \leftarrow T_0$
    \For{$k=1$ to $I_{max}$}
      \State $s' \leftarrow$ RandomNeighbor($s$)
      \State $\Delta = cost(s') - cost(s)$
      \If{$\Delta < 0$ or \text{rand()} < exp(-\Delta / T)}
        \State $s \leftarrow s'$
      \EndIf
      \State $T \leftarrow \gamma T$
    \EndFor
    \State \textbf{return} best observed $s$
  \EndProcedure
  \end{algorithmic}
  \end{algorithm}

  \section{Machine Learning Training and Evaluation Details}
  This appendix provides concrete hyperparameters, feature normalization steps,
  and evaluation pipelines used for k-NN and MLP models in `src/engine/mlRiskPipeline.ts`.

  \subsection{Feature Preprocessing}
  Features are normalized per-feature using training-set statistics (mean, std).
  Categorical features (victim severity) are one-hot encoded. Missing values are
  imputed with median (continuous) or mode (categorical).

  \subsection{k-NN}
  \begin{itemize}
    \item Distance metric: Manhattan (grid-oriented) and Euclidean (experimented).
    \item k selection: grid-search over $k \in \{1,3,5,7\}$ using cross-validation.
  \end{itemize}

  \subsection{MLP}
  \begin{itemize}
    \item Architecture: input\_dim - 32 - ReLU - 16 - ReLU - Sigmoid output.
    \item Loss: binary cross-entropy for survival label.
    \item Optimizer: Adam, learning rate $1e-3$.
    \item Batch size: 64, epochs: up to 200 with early stopping (patience 10).
  \end{itemize}

  \subsection{Evaluation Protocol}
  - Stratified 5-fold CV for robust estimates.
  - Report: Accuracy, Precision, Recall, F1, AUC-ROC and confusion matrix.
  - We additionally report calibration (Brier score) for $p_{survival}$ outputs.

  \section{Experimental Results Templates}
  Below are extended table templates and plotting suggestions to include in the
  final paper. Replace placeholders with real numbers from your runs.

  \begin{table*}[!t]
    \centering
    \caption{Detailed experimental results template (fill with your values)}
    \begin{tabular}{lrrrrrrrr}
      \toprule
      Run & VictimsSaved & AvgTime(s) & RiskExp & PathOptRatio & AmbUtil(\%) & Backtracks & ML-F1 & Notes \\
      \midrule
      A* (\alpha=1,\beta=0.5) & -- & -- & -- & -- & -- & -- & -- & time-prioritized \\
      Greedy (risk) & -- & -- & -- & -- & -- & -- & -- & risk-aware heuristic \\
      CSP (no heur) & -- & -- & -- & -- & -- & -- & -- & baseline CSP \\
      CSP (with heur) & -- & -- & -- & -- & -- & -- & -- & MRV+FC \\
      \bottomrule
    \end{tabular}
    \label{tab:results_template}
  \end{table*}

  Plot suggestions:
  \begin{itemize}
    \item Pareto front: VictimsSaved vs AvgRescueTime across weight sweep.
    \item Bar chart: Backtrack counts with/without CSP heuristics.
    \item Confusion matrix heatmap for ML models.
    \item Search trace time series: nodes expanded vs time for A* and Greedy.
  \end{itemize}

  \section{Example Decision Logs (Expanded)}
  Here are several example log entries demonstrating how to explain decisions.

  \begin{verbatim}
  2026-05-05T00:01:34Z | EVENT: road_blockage at (6,5) | AFFECTED: Amb1_trip2
  PREV_PLAN: Amb1: V1->M1 via P1 cost=120s risk=3.4
  NEW_PLAN: Amb1: V1->M1 via P2 cost=150s risk=0.6
  PRIORITIZED: time (alpha=1.0,beta=0.5) | JUSTIFICATION: Blocked edge on P1; alternate P2 increases time by 30s but reduces predicted risk exposure by 82%.

  2026-05-05T00:02:10Z | EVENT: new_victim at (9,9) | AFFECTED: CSP
  PREV_ASSIGN: Amb2 scheduled for V3,V5
  NEW_ASSIGN: Amb2 -> V3; Amb1 assigned V5 (rebalancing to respect capacity)
  PRIORITIZED: survival (ML suggests V_new p_survival=0.22)

  2026-05-05T00:03:50Z | EVENT: kit_depletion Amb1 | AFFECTED: subsequent pickups
  ACTION: Reallocate kit from base; postpone non-critical pickup
  JUSTIFICATION: Kits limited, prioritize critical victims
  \end{verbatim}

  \section{Writing and Page-Limit Tips}
  To fit the 6-page IEEE limit while preserving content quality:
  \begin{itemize}
    \item Use concise captions and combine related figures into subfigures.
    \item Move extended algorithm pseudocode and long tables to an appendix (as
      allowed by submission) or supplementary material repository.
    \item Use `\small` environment for long tables and `\scriptsize` for dense
      captions where necessary (but keep readability).
  \end{itemize}

  \section{References (suggested starter entries)}
  Add relevant citations to `refs.bib`. Example BibTeX entries:
  \begin{verbatim}
  @book{russell2016artificial,
    title={Artificial Intelligence: A Modern Approach},
    author={Russell, Stuart and Norvig, Peter},
    year={2016},
    publisher={Pearson}
  }

  @article{hart1968formal,
    title={A formal basis for the heuristic determination of minimum cost paths},
    author={Hart, P. E. and Nilsson, N. J. and Raphael, B.},
    journal={IEEE Transactions on Systems Science and Cybernetics},
    year={1968}
  }
  \end{verbatim}

  \section{Next Steps I can take for you}
  If you want, I can now:
  \begin{itemize}
    \item Insert your real group names and GitHub/LinkedIn links into `report.tex`.
    \item Generate a placeholder vector `figs/architecture.pdf` from a simple
      diagram description and commit it.
    \item Run the project's simulation harness (if available) to produce KPI CSVs
      and populate the result tables and figures.
  \end{itemize}

  Tell me which of these you'd like me to do next and provide any missing
  details (group names, GitHub repo URL, whether to run simulations).

  \end{document}
